% ═══════════════════════════════════════════════════════════════
% SLIDES - DS-03: Repaso general
% Klasse 9, Unidad 1+2 - Klausurvorbereitung
% ═══════════════════════════════════════════════════════════════

\documentclass[aspectratio=169, 11pt]{beamer}

\usepackage[utf8]{inputenc}
\usepackage[T1]{fontenc}
\usepackage[ngerman]{babel}

\usepackage{helvet}
\usepackage{palatino}
\newcommand{\seriftext}[1]{{\fontfamily{ppl}\selectfont #1}}

\usepackage{tikz}
\usetikzlibrary{calc}
\usepackage{tabularx}
\usepackage{booktabs}
\usepackage{array}

% ═══════════════════════════════════════════════════════════════
% FARBEN
% ═══════════════════════════════════════════════════════════════

\definecolor{spanischrot}{HTML}{C41E3A}
\definecolor{dunkel}{HTML}{222222}
\definecolor{mittelgrau}{HTML}{666666}
\definecolor{hellgrau}{HTML}{F5F5F5}

\definecolor{station1}{HTML}{4CAF50}
\definecolor{station2}{HTML}{2196F3}
\definecolor{station3}{HTML}{9C27B0}
\definecolor{station4}{HTML}{FF9800}
\definecolor{station5}{HTML}{E91E63}

% ═══════════════════════════════════════════════════════════════
% BEAMER-THEME
% ═══════════════════════════════════════════════════════════════

\usetheme{default}
\usecolortheme{default}

\setbeamercolor{structure}{fg=spanischrot}
\setbeamercolor{frametitle}{fg=dunkel, bg=white}
\setbeamercolor{title}{fg=white}
\setbeamercolor{subtitle}{fg=white}
\setbeamercolor{author}{fg=white}
\setbeamercolor{date}{fg=white}
\setbeamercolor{normal text}{fg=dunkel}

\setbeamerfont{frametitle}{size=\Large, series=\bfseries}
\setbeamerfont{title}{size=\LARGE, series=\bfseries}

\setbeamertemplate{navigation symbols}{}
\setbeamertemplate{footline}{%
    \hfill\textcolor{mittelgrau}{\insertframenumber}\hspace{1em}\vspace{0.5em}
}

% ═══════════════════════════════════════════════════════════════
% BEFEHLE
% ═══════════════════════════════════════════════════════════════

\newcommand{\es}[1]{\seriftext{\textbf{#1}}}
\newcommand{\de}[1]{{\small\textcolor{mittelgrau}{\textit{#1}}}}
\newcommand{\regel}[1]{%
    \tikz[baseline=(X.base)]\node[fill=spanischrot, text=white,
    font=\scriptsize\bfseries, inner sep=3pt, rounded corners=1pt](X){#1};%
}

% ═══════════════════════════════════════════════════════════════
% DOKUMENT
% ═══════════════════════════════════════════════════════════════

\begin{document}

% ═══════════════════════════════════════════════════════════════
% TITELFOLIE
% ═══════════════════════════════════════════════════════════════

{
\setbeamercolor{background canvas}{bg=spanischrot}
\begin{frame}[plain]
    \vfill
    \begin{center}
        {\usebeamerfont{title}\usebeamercolor[fg]{title}Repaso general}\\[8pt]
        {\usebeamercolor[fg]{subtitle}Klausurvorbereitung U1 + U2}\\[20pt]
        {\usebeamercolor[fg]{date}\small Spanisch · Klasse 9 · DS-03}
    \end{center}
    \vfill
\end{frame}
}

% ═══════════════════════════════════════════════════════════════
% ÜBERBLICK
% ═══════════════════════════════════════════════════════════════

\begin{frame}{Überblick: Was kommt in der Klausur?}
    \begin{columns}
        \begin{column}{0.5\textwidth}
            \textbf{Unidad 1:}
            \begin{itemize}
                \item Vokabeln (Alltag, Freizeit)
                \item Pretérito perfecto
            \end{itemize}

            \vspace{10pt}

            \textbf{Unidad 2:}
            \begin{itemize}
                \item Hausvokabeln + Möbel
                \item Bewegungsverben + Präp.
                \item Objektpronomen
            \end{itemize}
        \end{column}
        \begin{column}{0.5\textwidth}
            \begin{center}
                \fcolorbox{spanischrot}{hellgrau}{\parbox{0.85\textwidth}{\centering
                    \textbf{Klassenarbeit}\\[4pt]
                    19./20. Januar\\[4pt]
                    {\small Leseverstehen + Grammatik + Textproduktion}
                }}
            \end{center}
        \end{column}
    \end{columns}
\end{frame}

% ═══════════════════════════════════════════════════════════════
% STATIONENLERNEN ERKLÄRUNG
% ═══════════════════════════════════════════════════════════════

\begin{frame}{Heute: Stationenlernen}
    \begin{center}
        \textbf{5 Stationen -- du arbeitest selbstständig!}
    \end{center}

    \vspace{10pt}

    \begin{columns}
        \begin{column}{0.5\textwidth}
            \fcolorbox{station1}{white}{\parbox{0.9\textwidth}{\centering\small\textbf{1. Vocabulario}\\Vokabeln U1+U2}}

            \vspace{8pt}

            \fcolorbox{station2}{white}{\parbox{0.9\textwidth}{\centering\small\textbf{2. Pret. perfecto}\\Konjugation + Sätze}}

            \vspace{8pt}

            \fcolorbox{station3}{white}{\parbox{0.9\textwidth}{\centering\small\textbf{3. Verbos + Pron.}\\Bewegungsverben}}
        \end{column}
        \begin{column}{0.5\textwidth}
            \fcolorbox{station4}{white}{\parbox{0.9\textwidth}{\centering\small\textbf{4. Lesen}\\Comprensión lectora}}

            \vspace{8pt}

            \fcolorbox{station5}{white}{\parbox{0.9\textwidth}{\centering\small\textbf{5. Schreiben}\\Textproduktion}}

            \vspace{8pt}

            {\small\textit{Pro Station: ca. 14-17 Min}\\
            \textit{Lösungsblätter zur Selbstkontrolle!}}
        \end{column}
    \end{columns}
\end{frame}

% ═══════════════════════════════════════════════════════════════
% STATION 1: VOCABULARIO
% ═══════════════════════════════════════════════════════════════

{
\setbeamercolor{frametitle}{fg=white, bg=station1}
\begin{frame}{Station 1: Vocabulario U1 + U2}
    \begin{columns}
        \begin{column}{0.5\textwidth}
            \textbf{Hausräume:}

            \vspace{4pt}

            {\small
            \begin{tabular}{ll}
            \es{el salón} & Wohnzimmer \\
            \es{el dormitorio} & Schlafzimmer \\
            \es{la cocina} & Küche \\
            \es{el baño} & Badezimmer \\
            \es{el pasillo} & Flur \\
            \end{tabular}}
        \end{column}
        \begin{column}{0.5\textwidth}
            \textbf{Möbel:}

            \vspace{4pt}

            {\small
            \begin{tabular}{ll}
            \es{la cama} & Bett \\
            \es{el armario} & Schrank \\
            \es{el sofá} & Sofa \\
            \es{la nevera} & Kühlschrank \\
            \es{el televisor} & Fernseher \\
            \end{tabular}}
        \end{column}
    \end{columns}

    \vspace{10pt}

    \begin{center}
        \regel{Jetzt} $\rightarrow$ AB Aufgabe 1 (12 Items)
    \end{center}
\end{frame}
}

% ═══════════════════════════════════════════════════════════════
% STATION 2: PRETÉRITO PERFECTO
% ═══════════════════════════════════════════════════════════════

{
\setbeamercolor{frametitle}{fg=white, bg=station2}
\begin{frame}{Station 2: El pretérito perfecto}
    \begin{columns}
        \begin{column}{0.5\textwidth}
            \textbf{Konjugation \es{haber}:}

            \vspace{4pt}

            {\small
            \begin{tabular}{ll}
            yo & \es{he} \\
            tú & \es{has} \\
            él/ella & \es{ha} \\
            nosotros & \es{hemos} \\
            vosotros & \es{habéis} \\
            ellos & \es{han} \\
            \end{tabular}}
        \end{column}
        \begin{column}{0.5\textwidth}
            \textbf{Partizipien:}

            \vspace{4pt}

            {\small
            -ar → \textbf{-ado}\\
            -er/-ir → \textbf{-ido}

            \vspace{6pt}

            \textbf{Unregelmäßig:}\\
            poner → \textbf{puesto}\\
            hacer → \textbf{hecho}
            }
        \end{column}
    \end{columns}

    \vspace{10pt}

    \begin{center}
        \regel{Jetzt} $\rightarrow$ AB Aufgabe 2 (Konjugation + Sätze)
    \end{center}
\end{frame}
}

% ═══════════════════════════════════════════════════════════════
% STATION 3: VERBOS + PRONOMBRES
% ═══════════════════════════════════════════════════════════════

{
\setbeamercolor{frametitle}{fg=white, bg=station3}
\begin{frame}{Station 3: Verbos de movimiento + Pronombres}
    \begin{columns}
        \begin{column}{0.5\textwidth}
            \textbf{Bewegungsverben:}

            \vspace{4pt}

            {\small
            \begin{tabular}{ll}
            \es{bajar de} & herunterbringen \\
            \es{subir a} & hochbringen \\
            \es{poner en} & stellen \\
            \es{sacar de} & herausnehmen \\
            \es{llevar a} & hinbringen \\
            \es{traer de} & herbringen \\
            \end{tabular}}
        \end{column}
        \begin{column}{0.5\textwidth}
            \textbf{Objektpronomen:}

            \vspace{4pt}

            {\small
            \begin{tabular}{ll}
            \es{lo} & ihn \\
            \es{la} & sie (f) \\
            \es{los} & sie (m.Pl.) \\
            \es{las} & sie (f.Pl.) \\
            \end{tabular}

            \vspace{6pt}

            \textbf{Stellung:} VOR dem Verb!\\
            \es{La pongo en la mesa.}
            }
        \end{column}
    \end{columns}

    \vspace{10pt}

    \begin{center}
        \regel{Jetzt} $\rightarrow$ AB Aufgabe 3 (6 Sätze umformen)
    \end{center}
\end{frame}
}

% ═══════════════════════════════════════════════════════════════
% STATION 4: LESEVERSTEHEN
% ═══════════════════════════════════════════════════════════════

{
\setbeamercolor{frametitle}{fg=white, bg=station4}
\begin{frame}{Station 4: Comprensión lectora}
    \begin{center}
        \textbf{Strategie fürs Leseverstehen:}
    \end{center}

    \vspace{10pt}

    \begin{enumerate}
        \item \textbf{Fragen zuerst lesen} -- Was wird gefragt?
        \item \textbf{Text überfliegen} -- Worum geht es?
        \item \textbf{Schlüsselwörter markieren} -- Wo steht die Antwort?
        \item \textbf{Antworten formulieren} -- Mit eigenen Worten!
    \end{enumerate}

    \vspace{15pt}

    \begin{center}
        \regel{Jetzt} $\rightarrow$ AB Aufgabe 4 (Text + 5 Fragen)
    \end{center}
\end{frame}
}

% ═══════════════════════════════════════════════════════════════
% STATION 5: TEXTPRODUKTION
% ═══════════════════════════════════════════════════════════════

{
\setbeamercolor{frametitle}{fg=white, bg=station5}
\begin{frame}{Station 5: Producción de textos}
    \begin{center}
        \textbf{Strategie fürs Schreiben:}
    \end{center}

    \vspace{10pt}

    \begin{enumerate}
        \item \textbf{Erst planen} -- Was will ich schreiben?
        \item \textbf{Verschiedene Satzanfänge} -- Nicht immer ``Yo...''
        \item \textbf{Konnektoren nutzen} -- \es{también, pero, porque, y}
        \item \textbf{Korrekturlesen} -- Akzente, Genus, Verbformen
    \end{enumerate}

    \vspace{10pt}

    \textbf{Nützliche Strukturen:}\\
    \es{Mi habitación está... / En mi habitación hay... / Me gusta...}

    \vspace{10pt}

    \begin{center}
        \regel{Jetzt} $\rightarrow$ AB Aufgabe 5 (min. 5 Sätze)
    \end{center}
\end{frame}
}

% ═══════════════════════════════════════════════════════════════
% SELBSTEINSCHÄTZUNG
% ═══════════════════════════════════════════════════════════════

\begin{frame}{Selbsteinschätzung: Ampelkarte}
    \begin{center}
        \textbf{Wie sicher fühlst du dich?}
    \end{center}

    \vspace{10pt}

    \begin{columns}
        \begin{column}{0.33\textwidth}
            \begin{center}
                \fcolorbox{green!70!black}{green!20}{\parbox{0.8\textwidth}{\centering
                    \textbf{Grün}\\[4pt]
                    {\small Ich bin sicher!\\
                    Kann alles erklären.}
                }}
            \end{center}
        \end{column}
        \begin{column}{0.33\textwidth}
            \begin{center}
                \fcolorbox{yellow!70!black}{yellow!20}{\parbox{0.8\textwidth}{\centering
                    \textbf{Gelb}\\[4pt]
                    {\small Noch unsicher.\\
                    Muss üben.}
                }}
            \end{center}
        \end{column}
        \begin{column}{0.33\textwidth}
            \begin{center}
                \fcolorbox{red!70!black}{red!20}{\parbox{0.8\textwidth}{\centering
                    \textbf{Rot}\\[4pt]
                    {\small Brauche Hilfe!\\
                    Verstehe es nicht.}
                }}
            \end{center}
        \end{column}
    \end{columns}

    \vspace{15pt}

    \begin{center}
        {\small Bewerte dich für jedes Thema: Vokabeln, Pret. perf., Pronomen, Lesen, Schreiben}
    \end{center}
\end{frame}

% ═══════════════════════════════════════════════════════════════
% TIPPS FÜR DIE KLAUSUR
% ═══════════════════════════════════════════════════════════════

\begin{frame}{Tipps für die Klausur}
    \begin{columns}
        \begin{column}{0.5\textwidth}
            \textbf{Vor der Klausur:}
            \begin{itemize}
                \item Vokabeln mehrfach wiederholen
                \item Partizipien auswendig lernen
                \item Beispieltexte schreiben
                \item Genug schlafen!
            \end{itemize}
        \end{column}
        \begin{column}{0.5\textwidth}
            \textbf{In der Klausur:}
            \begin{itemize}
                \item Aufgaben erst durchlesen
                \item Zeit einteilen
                \item Bei Unsicherheit weitermachen
                \item Am Ende Korrekturlesen
            \end{itemize}
        \end{column}
    \end{columns}

    \vspace{15pt}

    \begin{center}
        \fcolorbox{spanischrot}{hellgrau}{\parbox{0.5\textwidth}{\centering
            \textbf{¡Buena suerte!}\\
            Viel Erfolg bei der Klausur!
        }}
    \end{center}
\end{frame}

% ═══════════════════════════════════════════════════════════════
% ABSCHLUSS
% ═══════════════════════════════════════════════════════════════

{
\setbeamercolor{background canvas}{bg=spanischrot}
\begin{frame}[plain]
    \vfill
    \begin{center}
        {\usebeamerfont{title}\usebeamercolor[fg]{title}¡Buen trabajo!}\\[20pt]
        {\usebeamercolor[fg]{subtitle}\small Fragen? Jetzt klären!}\\[10pt]
        {\usebeamercolor[fg]{subtitle}\small Klassenarbeit: 19./20. Januar}
    \end{center}
    \vfill
\end{frame}
}

\end{document}
